\chapter{Introduccción}

La primera vez que tuve un ud en mis manos, llevaba ya décadas tocando la guitarra y bastantes años haciendo lo propio con el violín. Mi bagaje musical me había llevado a estudiar con cierta profundidad formas musicales como el blues, el jazz, la música celta o el flamenco, siendo esta última a la que dedicaba la mayor parte de mi tiempo por aquel entonces. Sin embargo, mi exposición a la música árabe había sido practicamente nula. Todo lo más, la que pudiera llegarme por la influencia de esta en el flamenco y las cercanía que ambas música tiene en su sonoridad (elemento que, como supe más tarde, no es definitorio de la música árabe más que en lo superficial y para quien todo lo desconoce de ella, de la misma manera que sucede en el propio flamenco), amén de acaso alguna pieza de música arabo-andalusí.

No creo que mi circunstancia sea excepcional. Al contrario, creo que este mismo camino que me llevó hasta el ud es el que han seguido la mayor parte de sus interpretes occidentales, ya sean profesionales o simples aficionados. Es después de adquirir experiencia musical en otros instrumentos y estilos, y llevado por la curiosidad y el deseo de ampliar horizontes, como uno suele llegar a él si se es parte de una tradición cultural distinta a la de su origen.

El ud, instrumento rey de la música árabe como lo es el piano en la música clásica o la guitarra en la música popular occidental, es un instrumento minoritario y poco conocido en nuestro ámbito músical, de tal manera que es raro comenzar con él a adentrarse en el terreno de la interpretación musical, al contrario de lo que sí sucede en el mundo árabe.

Es con esta idea que escribo este libro, pensando en quienes comparten con mi yo de entonces **** para 

Recojo aquí no solo mi conocimiento del ud y lo que de él he recogido en este tiempo a partir de fuentes muy diversas, sino también mi experiencia tratando de aprender sus secretos y los de la música árabe desde esa circunstancia que, como digo, creo que será común a la mayor parte de los lectores de este libro.

****

El libro asume que el lector no tiene conocimiento previo del ud desde el punto de vista del intérprete. Tampoco se asume experiencia en la música árabe, tanto sea en lo referente a su teoría  y fundamentos como en la mera experiencia en calidad de oyente. Es decir, se considera que el lector es nuevo tanto en el ud como en el contexto musical y cultural en el que este se sitúa. Los elementos de ambas partes se detallaran a lo largo del libro, de tal manera que este puede servir no solo como una obra de referencia para músicos que deseen aprender a tocar el ud, sino también para quienes busquen adentrarse en la música árabe y su ámbito, bien como interpretes o como oyentes.

Por el contrario, sí que se asume un cierto conocimiento básico de teoría y notación musical. De no ser así, el lector deberá complementar alguno de los muchos textos disponibles sobre dichos temas, para después venir a este para profundizar en los aspectos particulares de la música árabe.

El contenido de este libro no está pensado a modo de método de aprendizaje, y no contiene ejercicios o estudios para ****